\section{Chapitre~\ref{sec:synthesis}. Toute la machine}

Ce chapitre de synthèse n'apporte pas d'expériences nouvelles: chaque ligne s'appuie sur le chapitre où l'affirmation est examinée et sur les mêmes sources; le bus \nt{GLUT} et la commande du flux par \nt{GABA} sont solidement établis, tandis que les rôles des canaux de diffusion et leur fonctionnement conjoint sont pour l'essentiel une hypothèse.

{\footnotesize
\setlength{\tabcolsep}{3pt}\renewcommand{\arraystretch}{1.15}
\begin{longtable}{>{\raggedright}p{0.5cm}>{\raggedright}p{4.0cm}>{\raggedright}p{5.6cm}>{\raggedright}p{2.3cm}>{\raggedright\arraybackslash}p{1.7cm}}
\toprule
N\textsuperscript{o} & Affirmation & Expérience et ce qui est mesuré & Source & Statut \\
\midrule\endfirsthead
\toprule
N\textsuperscript{o} & Affirmation & Expérience et ce qui est mesuré & Source & Statut \\
\midrule\endhead
1 & Pour $8.6\cdot10^{10}$ neurones, il n'existe que quatre sortes de signaux de commande~\eqref{eq:machine} & Comptage des noyaux cellulaires dans un homogénat de cerveau (fractionneur isotrope): $86.1\pm8.1$ milliards de neurones chez l'homme adulte & \cite{Azevedo2009} & mesuré \\
2 & Proposition~\ref{prop:architecture}, deux premières couches: les données circulent sur le bus d'adresses \nt{GLUT}, le signe de la connexion est fixé par l'émetteur, les neurones \nt{GABA} orientent et normalisent le flux (chapitres~\ref{sec:neurontypes}, \ref{sec:gaba}) & Un seul neurotransmetteur principal par neurone (revue des mesures); l'inhibition directe rétrécit la fenêtre dans laquelle le neurone pyramidal additionne ses entrées; la division par l'activité totale est mesurée dans de nombreuses régions & \cite{Strata1999,Pouille2001,Carandini2012} & mesuré \\
3 & Les éléments ne sont pas fiables: une synapse perd la plupart des impulsions (tabl.~\ref{tab:computer}, chapitre~\ref{sec:code}) & Tranches d'hippocampe, stimulation minimale: plus de la moitié des impulsions ne provoquent aucune réponse dans CA1; les échecs du seuil et de la conduction axonale sont exclus, la cause est la libération probabiliste du neurotransmetteur & \cite{AllenStevens1994} & mesuré \\
4 & Le cerveau fonctionne avec $\sim20$ watts, en moyenne environ une impulsion par seconde et par neurone (chapitre~\ref{sec:code}); les ingénieurs n'ont pas encore construit une telle machine & Estimation~\eqref{eq:energy} à partir des coûts mesurés: environ $10^9$ molécules d'ATP par impulsion, travail des synapses compris (cortex de rongeur); une estimation détaillée pour le cortex donne une fréquence encore plus basse. État de la technique d'après une revue & \cite{Attwell2001,Lennie2003,MehonicKenyon2022} & théorie \\
5 & L'apprentissage de la plupart des synapses est le produit d'une trace locale et d'un seul nombre commun $\delta$ (deuxième équation de~\eqref{eq:machine}, chapitre~\ref{sec:dofa}) & Les neurones \nt{DOPA} du singe répondent à l'erreur de prédiction de la récompense; dans des tranches de striatum, la dopamine n'agrandit une épine que si elle arrive $0.3$--$2$\,s après l'entrée glutamatergique: la trace est mesurée & \cite{Schultz1997,Yagishita2014} & mesuré \\
6 & Un fil d'erreur propre à chaque sortie n'existe que dans le circuit d'apprentissage moteur (chapitre~\ref{sec:motorlearning}) & Cervelet: la coïncidence du signal de la fibre grimpante avec une entrée affaiblit cette entrée; chez le singe, les impulsions complexes des cellules de Purkinje portent la direction de l'erreur de mouvement et déclenchent une correction & \cite{Ito2001,Herzfeld2018} & mesuré \\
7 & Chaque canal de diffusion est un faisceau de quelques lignes (proposition~\ref{prop:bundle}) & Pour \nt{DOPA}: dans une même tâche, des neurones différents portent la récompense, le mouvement et les propriétés du stimulus; l'erreur rapide et le niveau lent de dopamine sont distincts. Pour \nt{SERO}, \nt{NORA}, \nt{ACET}, cette décomposition est moins étudiée & \cite{Engelhard2019,Hamid2016,Mohebi2019} & mesuré \\
8 & \nt{SERO}: régulateur de niveau et, par hypothèse, horizon $\tau_\gamma$ (chapitre~\ref{sec:serocomputer}) & Auto-inhibition: les agonistes de leurs propres récepteurs 5-HT$_{1A}$ inhibent les neurones sérotoninergiques du raphé (mesuré). L'horizon est proposé en théorie; l'expérience la plus forte: l'activation optogénétique de ces neurones chez la souris allonge l'attente d'une récompense différée & \cite{Sprouse1987,Doya2002,Miyazaki2014} & hypothèse \\
9 & \nt{NORA}: le gain $g$ choisit entre chercher et décider (chapitre~\ref{sec:nora}) & Hausse du rapport signal/bruit: chez le singe éveillé, la noradrénaline dans le cortex auditif supprime l'activité de fond plus fortement que la réponse aux vocalisations des congénères (mesuré); le gain multiplicatif est un modèle; le rôle dans le choix entre chercher et décider est une théorie & \cite{Foote1975NE,ServanSchreiber1990,AstonJones2005} & hypothèse \\
10 & \nt{ACET}: bascule entre écriture et lecture (chapitre~\ref{sec:acet}) & Les trois actions (affaiblissement des connexions internes, renforcement de l'entrée, facilitation de la plasticité) sont mesurées; la bascule des modes est soutenue indirectement par une expérience à la scopolamine chez l'homme & \cite{HasselmoBower1992,Gil1997,Atri2004} & hypothèse \\
11 & La formule~\eqref{eq:machine} décrit la machine entière & C'est un résumé des modèles des chapitres~\ref{sec:glutcomputer}--\ref{sec:acet}, chaque partie s'appuie sur son chapitre; l'ensemble n'a pas été vérifié par l'expérience & déduction du chapitre & hypothèse \\
12 & Les boutons fonctionnent de concert sans commande centrale: erreur, surprise, écriture, retour (fig.~\ref{fig:timeline}) & Pas d'expérience: le schéma est qualitatif. Maillons isolés: théorie de la division du travail entre \nt{NORA} et \nt{ACET}, et lien entre la pupille et la vitesse d'apprentissage chez l'homme & \cite{YuDayan2005,Nassar2012} & hypothèse \\
13 & L'apprentissage par un seul signal commun est d'autant plus lent que plus de neurones influent sur le résultat (proposition~\ref{prop:broadcastcost}) & Calcul des courbes d'apprentissage d'un réseau linéaire: l'apprentissage par perturbation des sorties est plus lent que la descente de gradient d'un facteur égal au nombre de sorties & \cite{Werfel2005} & théorie \\
14 & On ignore comment le cortex règle ses circuits multicouches; candidats: bouffées d'impulsions, calculs dans les branches du neurone, apprentissage auto-supervisé par prédiction & Bouffées comme porteuses de l'erreur: modèle; seul un réseau de $5$--$8$ couches a su reproduire la réponse d'un modèle détaillé de neurone cortical: modèle; potentiels calciques dans les branches des neurones corticaux humains, réalisant le OU exclusif: mesuré sur tranches; apprentissage auto-supervisé: revue des indices & \cite{Payeur2021,Beniaguev2021,Gidon2020,Keller2018} & hypothèse \\
\bottomrule
\end{longtable}}
